\documentclass[a4paper,10pt,onecolumn]{article}

% Wider margins (compared to default)
%\usepackage[margin=2.5cm]{geometry}
\usepackage[left=2.5cm,right=2.5cm,top=2.5cm,bottom=2.5cm]{geometry}


% English support (typography and hyphenation)
\usepackage[american]{babel}
\usepackage{csquotes}

% For math
\usepackage{amsmath}

% Unicode encoding
\usepackage[utf8]{inputenc}

% Color box
\usepackage{tikz}
\usepackage[most]{tcolorbox}


% Better default font (Libertine and Inconsolata)
\usepackage[ttscale=.875]{libertine}
\usepackage[scaled=0.96]{zi4}


% Biber
\usepackage[backend=biber,
            style=apa,
            % backref=true,
            % backrefstyle=three,
            maxcitenames=3,
            maxbibnames=99,
            apamaxprtauth=99,
            natbib=true]{biblatex}
\DeclareLanguageMapping{american}{american-apa}

% Figure & graphics
\usepackage{tikz}
\usepackage{wrapfig}
\usepackage{graphicx}

% Figure caption
\usepackage[labelsep=period]{caption}
\captionsetup[figure]{font=small,skip=0pt}
\renewcommand{\captionfont}{\small\sffamily}
\renewcommand{\captionlabelfont}{\small\sffamily\bfseries}

% Hyperref
\usepackage{xcolor}
\definecolor{blendedblue}{rgb}{0.2, 0.2, 0.6}
\definecolor{blendedred}{rgb}{0.8, 0.2, 0.2}
\usepackage[bookmarks=true,
            breaklinks=true,
            pdfborder={0 0 0},
            citecolor=blendedblue,
            colorlinks=true,
            linkcolor=blendedblue,
            urlcolor=blendedblue,
            citecolor=blendedblue,
            linktocpage=false,
            hyperindex=true,
            linkbordercolor=white]{hyperref}
\usepackage{hyperref}
\hypersetup{colorlinks=true}


% Bibliography file
\bibliography{article.bib}

% Headers
\usepackage{fancyhdr}
\pagestyle{fancy}
% \fancyhf{}
\rhead{\footnotesize \sf \today}
\lhead{\footnotesize \sf Rougier \& Timmer \textbullet~Ten Simple Rules for Scientific Fraud \& Misconduct}
\rfoot{\footnotesize \sf \thepage}
\cfoot{}
\lfoot{}

\makeatletter
\renewcommand{\maketitle}{\bgroup\setlength{\parindent}{0pt}
\begin{flushleft}
  \textbf{\huge\@title\\}
  \vspace{5mm}
  \@author
\end{flushleft}\egroup
}
\makeatother

\title{\Large Ten Simple Rules for Scientific Fraud \& Misconduct}
\author{%
  \textbf{Nicolas P. Rougier}$^{1,2,3,*}$ \textbf{and John Timmer}$^{4}$\\
  \begin{footnotesize}
    $^{1}$INRIA Bordeaux Sud-Ouest
          Talence, France
    $^{2}$Institut des Maladies Neurodégénératives,
          Université de Bordeaux, CNRS UMR 5293,
          Bordeaux, France
    $^{3}$LaBRI, Université de Bordeaux, Institut Polytechnique de Bordeaux,
          CNRS, UMR 5800, Talence, France
    $^{4}$Ars Technica, Condé Nast, New York, NY, USA\\
    $^{*}$Corresponding author:
          \href{mailto:Nicolas.Rougier@inria.fr}{Nicolas.Rougier@inria.fr}
  \end{footnotesize}
}

\date{\today}

\begin{document}
\maketitle

\begin{tikzpicture}[remember picture, overlay]
    \node[xshift=-2.85cm,yshift=-2.85cm] at (current page.north east) {
        \includegraphics[width=3cm]{stample}
    };
\end{tikzpicture}

%% \begin{figure}[h]
%%   \begin{center}
%%     \noindent \includegraphics[width=\textwidth]{fraud-kit}
%%   \end{center}
%%   \caption{My First Fraud Kit. Copyright (c) 2012 Aurich Lawson. Ars Technica
%%     \citep{timmer:2012}.}
%% \end{figure}


\begin{tcolorbox}[breakable, pad at break*=1mm,
    colback=yellow!5, arc=0pt, outer arc=0pt, boxrule=0pt, frame hidden]
\begin{small}
  \noindent {\sf \textbf{Disclaimer.} We obviously do not encourage scientific
    fraud nor misconduct. The goal of this article is to alert the reader to
    problems that have arisen in part due to the Publish or Perish imperative,
    which has driven a number of researchers to cross the Rubicon without the
    full appreciation of the consequences. Choosing fraud will hurt science,
    end careers, and could have impacts on life outside of the lab. If you're
    tempted (even slightly) to beautify your results, keep in mind that the
    benefits are probably not worth the risks.}\par
\end{small}
\end{tcolorbox}

\section*{Introduction}

So, here we are! You've decided to join the dark side of Science. That's great!
You'll soon discover a brand new world of surprising results,
non-replicable experiments, fabricated data, and funny statistics. But it's not without risks: fame and shame,
retractions and lost grants, and... possibly jail. But you've made your choice, so now you need to know how to manage these risks. Only a few years ago, fraud and misconduct was a piece of
cake (See the \href{https://en.wikipedia.org/wiki/The_Turk}{Mechanical Turk},
\href{https://en.wikipedia.org/wiki/Charles_Redheffer}{Perpetual motion
  machine}, \href{https://en.wikipedia.org/wiki/Great_Moon_Hoax}{Life on Moon},
\href{https://en.wikipedia.org/wiki/Piltdown_Man}{Piltdown man},
\href{https://en.wikipedia.org/wiki/Water_memory}{Water memory}). But there are lots of new players in town (\href{https://pubpeer.com}{PubPeer},
\href{http://retractionwatch.com}{RetractionWatch},
\href{https://forbetterscience.com}{For Better Science},
\href{http://blogs.discovermagazine.com/neuroskeptic/}{Neuroskeptic} to name
just a few) who have gotten pretty good at spotting and reporting fraudsters.
Furthermore, publishers have started to arm themselves with high-tech tools, and your fellow scientists are willing to name and shame you \href{https://www.youtube.com/watch?v=FXaOqwanWnU/}{on social media}. To
commit fraud or misconduct without getting caught in 2017 is a real
challenge and requires serious dedication to your task. While you'll never be smarter than an entire community of scientists, we're generously giving you some
simple rules to follow in your brand new career in fraud (see also \citep{timmer:2012}
for a set of complementary rules). Of course, neither results or (lack of) consequences are 
guaranteed.



\section*{Rule 1: Misrepresent, falsify, or fabricate your data}

In order to start your life as a scientific fraudster, the
first thing you need to do is learn how to convincingly
misrepresent, falsify, or fabricate data. If you're still hesitant about
embracing the dark side of science, you can start with a slight
misrepresentation to support your hypothesis --- a hypothesis you're sure is right anyway. Some classic techniques are
described in more detail below (see the seventh rule of \citep{rougier:2014}) but there are a number of others. If your data include numbers, there are \href{https://en.wikipedia.org/wiki/Misleading_graph}{many
  techniques} for making a graph misleading \citep{wainer:1984,raschke:2008}. If you have images, you
can change your results using
software such as Photoshop or GIMP, which allow you to change pretty
much anything in a figure \citep{cromey:2012, hendricks:2011}. But be careful, as an ever-growing number of
journals are now using forensic software to spot image manipulation
\citep{white:2007, rossner:2004,hsu:2015} (see also
\href{http://journals.plos.org/ploscompbiol/s/figures#loc-image-manipulation}{PLoS
  image manipulation recommendation} and \href{http://jcb.rupress.org/content/166/1/11}{Rockefeller University Press' discussion of the issue}). You can even test your whether your image has issues that could tip off an editor using the online tool \href{https://29a.ch/photo-forensics/#forensic-magnifier}{Forensically}
by Jonas Wagner. This might dissuade you from manipulating your image, but don't worry --- there are plenty of other options available to you.\\


None of the previous techniques will change the statistics of your
data, so it might be good to consider your options here. Starting with real
data, you only need to change a few points in order to take a non-significant
result and turn it into something with an astoundingly highly significance. Just see what's possible using the
\href{http://shinyapps.org/apps/p-hacker/}{p-hacker} application
\citep{schonbrodt:2015} or the interactive applet from
\citep{aschwanden:2015}. The advantage of tweaking real data is that the results  look both good and not very
suspicious. The disadvantage is that it requires some real data in the first
place, which means carrying out some actual experiments. Who wants that? Instead, you can fabricate an entire data set using software. With just a little programming skill, you can choose a correlation coefficient and X and Y mean/SD, and generate a data set with these precise
statistics \citep{matejka:2017}. (Do not pick one of the datasaurus sets because if you
plot it, it might attract the attention of an alert reader.) Whatever option you choose, make sure to
have a backup story in case people start asking about the details of the experiments. A
number of misconduct cases have been detected with just a few questions
\citep{vastag:2006}.



\section*{Rule 2: Hack your results}

If you are reluctant to manipulate your data, you still have the option of
searching through your results to find anything that reaches significance 
(a.k.a. p-hacking). This can provide an appealing alternative to scientific misconduct. What is the p value of your NHST (Null Hypothesis
Significance Test)? If it's close to your field's standard one of the many qualification
proposed by
\href{https://mchankins.wordpress.com/2013/04/21/still-not-significant-2/}{Matthew
  Hankins} to report your results? Can't you use expression such as {\em nearly
  acceptable level of significance (p=0.06)}, {\em very closely brushed the
  limit of statistical significance (p=0.051)} or {\em weakly statistically
  significant (p=0.0557)}? While these statements don't make much sense, they might
be sufficient to convince a naive reviewer or reader. 


If you're
suffering from a bi-polar p-value disorder, however (a term coined by
\href{http://daniellakens.blogspot.fr/2014/05/the-probability-of-p-values-as-function.html}{Daniel
  Lakens}), you'll want to reach the magical 0.05 threshold. It might be much easier than you think to portray your data as significant! The correct method for calculating the
infamous p-value is barely understood by many students, teachers, and
researchers \citep{haller:2002,lecoutre:2003}. The current situation is so bad that the American Statistical Association had to release a
statement to explain p's correct usage \citep{wasserstein:2016,};
Nature has issued a warning as well
\citep{baker:2016}. Consequently, before hacking your results, it might be wise
to re-read your classics
\citep{simmons:2011,cumming:2012a,cumming:2012b,colquhoun:2014}. A warning for those of you in psychology: you have to take extra-precaution because this field has drawn attention for the unusual distribution of
p-values smaller than 0.05 \citep{hartgerink:2016,bakker:2012}. There's now software called statcheck that can be used to recalculate the p values in your paper \citep{nuijten:2015,epskamp:2016}, making it difficult to get away with misleading statistics. Although some have called checking the results of others {\em
  methodological terrorism} \citep{finske:2016}, it might be enough to scare some away from the dark side of Science.


\section*{Rule 3: Copy/paste from others}

It's not enough to have (fake) results; you still need to publish them if you want to gain fame (and tenure, and grants). Writing is a tedious task
and can be a fair amount of work. Summarizing the state of the art in your
field will force you to actually read about what your colleagues have been doing
over the past few years. It is a very time consuming task. But in doing that reading, you may find that one of these
colleagues wrote a nice introduction to the field or a wonderful summary of its current state. If so, why bother writing a new one? It's much simpler to copy/paste what he/she has written.  Plagiarism is the nuts and bolts of scientific misconduct
\citep{neuroskeptic:2017}, be it literal copying, substantial copying or
paraphrasing (see
\href{https://www.elsevier.com/editors/perk/plagiarism-complaints}{definitions
  from Elsevier} and the committee on publication ethics' (COPE) procedure for
handle plagiarism in a
\href{https://www.elsevier.com/__data/assets/pdf_file/0005/72815/plagiarism-A.pdf}{submitted}
or
\href{https://www.elsevier.com/__data/assets/pdf_file/0020/72830/plagiarism-B_0.pdf}{published}
article). 

Of course, you cannot literally take an excerpt from an
article to insert it in yours. You'd run the risk of creating a situation where the reviewer found his own text in yours. Consequently, if you
really intend to plagiarize, you'll need to transform the text drastically to
avoid detection by humans \citep{dorigo:2015,} and more importantly, to avoid
detection by dedicated software (e.g. iThenticate). This software is used by
universities and journals for tracking people just like you. As the databases behind these services grow, plagiarism is becoming harder and harder. Even
self-plagiarism is not allowed by a vast majority of scientific journals. All of this means that you have to be really innovative in your plagiarism techniques \citep{long:2009}. We've reached the point where writing your own text might be an easier option.

\section*{Rule 4: Write your own peer-review}

Even if you avoid all the pitfalls discussed above, you're in the clear yet. Depending on the
journal you target, you may still encounter problems with the review
process. If reviewers are picky, they may ask annoying questions,
request more information or experiments, ask for a major revision and/or reject your submission entirely. Can you imagine how this would make you feel after all your hard work
fabricating the data and plagiarized the article? There's even the risk that they'll notice your malpractice. Fortunately, there is a simple solution: write your
own review! 

It is surprisingly easy to do. As you submit, you will often be asked to give name of possible reviewers. Just
provide phony names, along with email addresses that will be redirected to
your mailbox. You will soon receive an invitation to review your own work, and you're then
free to state how brilliant your own work is. Of
course, you'll have to write a review that looks like an actual review. If
you're Machiavellian, you can introduce some factual errors in your
manuscript, then report them in your review, making it look thorough. Make sure not to send you review before the deadline because as
\href{https://arstechnica.com/science/2017/04/107-cancer-papers-retracted-due-to-peer-review-fraud/}{reported
  by Elizabeth Wager} \citep{stigbrand:2017}, editors have
spotted fake reviews in part because reviewers responded promptly. Unfortunately, editors and
publishers are now aware of the scam \citep{ferguson:2014} and  have
taken counter-measures \citep{haug:2015}. For example, some of them no longer offer authors the option of recommending reviewers or if they do, the
recommendation is restricted to a list of certified reviewers. If you insist on
writing your own peer-review, you'll have to be creative and find new ways to
game the system.



\section*{Rule 5: Take advantage of predatory publishers}

If you're worried that peer review will reveal your misconduct, you still have opportunities for publishing your results. There are many predatory publishers on
the internet \citep{shen:2015}. These predators will
publish just anything
(see \href{https://en.wikipedia.org/wiki/International_Journal_of_Advanced_Computer_Technology#Publication_controversy}{``Get
  Me Off Your Fucking Mailing List''} in the International Journal of Advanced
Computer Technology (2005, 2014)) and you have a 100\% chance of publication with a lighting fast review - less than 24h for some
journals. Finding a predatory publisher used to be easy thanks to a list
created and maintained by Jeffrey Beal. Although not all researchers (and publishers)
agreed with his categorization, this list provided an incredibly useful resource. Unfortunately, Beal's list was shut down just a few months ago
\citep{straumsheim:2017} and the
soon-to-be released replacement one will not be free \citep{silver:2017}. Still,
you can take advantage of the
\href{http://thinkchecksubmit.org}{Think/Check/Submit} website, which provides a
easy-to-use checklist that researchers can refer to when they are investigating
whether a journal is trustworthy. You'll obviously just want take the opposite of their sage recommendations --- like making sure \textbf{not} to pick a journal from the list of the 
\href{https://oaspa.org}{Open Access Scholarly Publishers Association} since
these journals have been thoroughly reviewed and can be trusted. Other anti-tips include finding a journal that none of your colleagues know about, looking for an editorial board that's made of
unknown or non-existent people \citep{sorokowski:2017} or dogs \citep{kennedy:2017}, and finding a publisher without any address / email / telephone number. Of course, you
will have to pay exorbitant fees to have your work published, but this is the
price to pay to have an expedited and complacent review, assuming there's any review at
all. Before making you decision where to publish, make sure to check
for the retraction fees, which may be much higher than the initial publishing
fees in some cases. But who wants to retract anyway?

\section*{Rule 6: Don't give access to your code and data}

You managed to have your article published, congratulations! You'll probably
feel a bit of a buzz, and your astonishing conclusion may draw some
newspaper headlines if your research is sexy enough. Don't get too excited about this, though. Attention means an increased chance that other researchers will start asking to access your
material, assuming the publisher has not already requested it prior to
publication. Of course, that's very bad news for you if your data
has been fabricated or if your statistical analysis is not really academic. You definitely cannot give others access to your raw if it doesn't exist! Fortunately, researchers have been avoiding sharing their data for decades, and the same old
reasons still apply
\citep{roche:2014}: ``Oh you know, my data cannot be anonymized''; ``You
would not understand the structure of my data''; ``A lot of money has been
invested, I cannot give it for free''; etc. You can even tell anyone
asking for your data that they are {\em a research parasite} \citep{longo:2016}
although that approach seems to have drawn some unwanted attention
(\url{http://researchparasite.com}).  Technological excuses might also be worth
a try. When \citep{collberg:2014,collberg:2015} tried to access the code used in
several publications, they get some really great answers: ``Too busy to help'';
``Code will be available soon''; ``Bad backup practice''; etc. Just pick one
you like. With the increasing pressure by funding bodies (NIH, NSF,
Europe, etc) to share data, it might be good to pretend (and only pretend) to be a
supporter of Open Data by inserting a simple ``data available upon request'' or
``data is available from my website''. This does not come close to meeting the
FAIR principles (findable, accessible, inter-operable, re-usable,
\cite{wilkinson:2016}), but who cares if you can give the illusion you're not
afraid of sharing your data? Don't get us wrong, though - you will never answer a
request, and the data will never appear on your website.


\section*{Rule 7: Do not allow for replication outside your lab}

Here comes the tricky part: replication. It may surprise you, but some researchers
may want to replicate your results using the methods explained in your
article. \citep{nosek:2015} replicated 100 experiments reported in three high-ranking psychology journal. Or tried, at least; they
concluded that some experiments were not replicable at all. There are even scientific journals that are now dedicated to replications
(e.g. \url{http://rescience.github.io}). 

If people try to replicate your work
and do not get the same results, you have a problem. They may insist on seeing your actual
data and, if you refuse, you might be suspected of
fraud or misconduct. You can try rule 6, but if there are several failed replications, it will be much
harder to deny access to your data. One easy way to avoid this sort of disaster is
to claim in your article that the experiment/study has been already
replicated in your lab. The subliminal message to the reader is ``Don't bother
to replicate this, we've already taken care of that and it works, believe us. Don't
waste your time, you will thank us later''. If this doesn't work, you can also
try issuing a {\em Do Not Replicate order} as proposed by
\href{https://mchankins.wordpress.com/2014/07/12/do-not-replicate/}{Matthew
  Hankins} (who is kind enough as to provide a template). It is not yet widely
used , but with the replication crisis we're facing, we're confident it will soon become popular.

\section*{Rule 8: Never, ever, retract your results}

If you've made a genuine (and big) mistake in your work, there is no problem in
asking for the retraction of your paper  \citep{miller:2006}. It's a
behavior that is actually rewarded according to \citep{lu:2013} (positive citation benefits for prior work), even though in some cases, it can take several years \citep{trivers:2009}. However, if you've been
engaged in fraud, having your paper retracted is like an admission of guilt. It's something to be avoided. Retraction could start with a simple comment to the
editor saying there may be a problem with you paper. If you're lucky, the comment will never be published
\citep{trebino:2009}. Otherwise, you'll have to give a convincing answer, or the editor could decide to retract your paper
without your consent. So it's critical to act quickly, defusing the crisis with a simple corrigendum admitting a bad - but not fatal - error during preparation of the publication.  Don't hesitate to publish as many corrigendum as
necessary to make critics happy. You can drag this out for several years, which is hopefully enough time for people to forget about the issues.

\section*{Rule 9: Don't get caught. Deny if caught.}

If you intend to persist in a rogue scientific career, you have to be aware that
you're likely to get caught sooner or later. The number of researcher to have never been caught is probably low; only the best can hope to be caught only after their death
\citep{degroote:2016,}. But being caught while you're still breathing does not mean you career will come to
an end. There is a set of simple rules to follow if you need to deny scientific misconduct:
\begin{itemize}
  \item If you're first author, explain you were supervised by the last author and
    had no choice.
  \item If you're last author, explain you were not aware of the misconduct of
    the first author.
  \item If your name is not first nor last, claim that you didn't even know your
    name appeared in the publication.
  \item If your name appears because of {\em gift authorship}, just admit it.
    It might be considered misconduct, but fraud is much worse.
  \item The intern, who cannot be contacted because he/she left academia, is
    responsible for everything
  \item Send threatening letters to those who have spotted your misconduct
  \item Follow through on those threats and sue'em all
\end{itemize}
If you're lucky enough, one of the above will work and things should be back to
normal. All of this is easier, of course, if you're supported by you employer because you've been
labeled a rising star.

\section*{Rule 10: Be creative (for once)}

If you look at the annual list of
\href{http://www.the-scientist.com/?articles.view/articleNo/47813/title/Top-10-Retractions-of-2016/}{top
  ten retractions} compiled by Adam Marcus and Ivan Oransky (co-founders of the
RetractionWatch), you'll realize that all the tactics mentioned above are
already quite well known by the research community. If you want to stay off the radar
while committing fraud and misconduct, you'll want to be creative and invent
your own rules. 

\section*{Conclusion}

By following the simple rules above, you should get scientific glory if only temporarily. The downside is that it could be followed by jail time \citep{grant:2015}. A former researcher has been sentenced to
57 months jail and to pay-back 7.2 millions dollars. Science has been and is
still poisoned by fraud and misconduct, but it is now fighting back with increasingly high-tech tools \citep{buranyi:2017}. Today, the risks that come when you engage in
fraud and misconduct are really high, and the chances of being
caught have gone up. So you'd better think twice before committing misconduct, or
your name will soon appear in Wikipedia's
\href{https://en.wikipedia.org/wiki/Scientific_misconduct#Notable_individual_cases}{hall of shame}.


% -------------------------------------------------------------- References ---
\renewcommand*{\bibfont}{\small}
\printbibliography[title=References]
\end{document}

